\chapter{系统安全设计}

\section{安全扩展的引入背景}

本文前面几章主要讨论的是细粒度分类性能,但从系统实现角度看,仅关注干净样本上的准确率并不充分。视觉模型在推理阶段可能受到对抗性扰动影响,即输入图像只经过很小的像素级修改,模型预测就会发生明显偏移。对于基于 CLIP 的提示学习系统而言,如果推理阶段完全不考虑这一风险,那么模型在实际部署时仍可能因为输入异常而产生不稳定输出。

基于这一考虑,本项目在原有动态提示分类系统之外,增加了一个独立的安全扩展模块。该模块并不改变训练主流程,而是在推理时对输入图像进行检测和防御,尽量降低对抗性扰动对特征表示和最终分类结果的影响。相关设计和实现主要记录在项目根目录的 \texttt{SECURITY.md}、\texttt{models/adversarial\_defense.py} 和 \texttt{models/robust\_custom\_clip.py} 中。

\section{安全机制原理}

项目采用的核心思路来自 Test-time Counterattack (TTC),即在测试阶段利用模型自身的视觉编码器对输入图像进行反向扰动和特征修正\cite{xing2025ttc}。与对抗训练不同,TTC 不要求在训练阶段额外加入大规模对抗样本,而是把主要计算放在推理阶段完成。

其基本思想可以概括为两步。第一步是检测: 对输入图像施加轻量扰动,比较扰动前后在 CLIP 视觉特征空间中的差异。如果特征变化异常大,则说明该输入可能包含对抗性成分。第二步是防御: 在检测结果基础上对图像执行反向攻击式修正,让修正后的图像特征尽量回到更稳定的区域,再送入后续分类流程。

项目文档中把这一关键量记为特征差异比率 \(\tau\),其定义为
\begin{equation}
\tau = \frac{\lVert f_{att} - f_{orig} \rVert}{\lVert f_{orig} \rVert},
\end{equation}
其中 \(f_{orig}\) 为原始图像的视觉特征,\(f_{att}\) 为加入扰动后的视觉特征。若 \(\tau\) 超过设定阈值,系统会认为当前输入更可能受到攻击,从而采用更强的防御策略。

\section{安全模块结构}

\subsection{TestTimeCounterattack}

\texttt{models/adversarial\_defense.py} 中的 \texttt{TestTimeCounterattack} 是安全模块的核心实现。它的主要职责是在推理阶段生成反向扰动,并根据特征差异自适应调整多步扰动的权重。

具体来说,该模块先对输入图像初始化一个有界扰动 \(\delta\),然后通过多步更新减小扰动图像与原始图像在特征空间中的偏差。项目实现中支持 \(\ell_{\infty}\) 与 \(\ell_2\) 两类扰动约束,默认配置使用 \(\ell_{\infty}\) 范数。每一步产生的扰动都会被记录,最后根据 \(\tau\) 的大小决定使用单步轻量修正还是多步加权修正。

\subsection{AdversarialDetector}

\texttt{AdversarialDetector} 负责对输入图像进行轻量检测。其实现方式比完整防御更简洁: 先给图像加入小范围随机扰动,再比较扰动前后的视觉特征差异。如果差异超过阈值,则把当前图像标记为潜在对抗样本。这样做的好处是开销相对较小,适合在自动模式中作为前置筛选。

\subsection{RobustCustomCLIP}

\texttt{models/robust\_custom\_clip.py} 将防御系统与原有的动态提示模型集成起来,形成 \texttt{RobustCustomCLIP} 和 \texttt{RobustCustomCLIPCoCoOp} 两个鲁棒模型接口。它们保留了提示学习与分类的主体流程,同时在图像进入视觉编码器之前插入防御步骤。

从代码结构看,安全功能被设计成一个推理时可选模块,而不是直接改写原有训练器。这样做有两个实际好处。第一,原始分类模型和鲁棒模型可以共用大部分提示学习实现,减少重复代码。第二,系统可以按需切换安全模式,不会影响正常训练流程。

\section{推理流程与模式设计}

集成安全模块后的推理流程可概括为:

\begin{enumerate}
\item 输入图像首先进入防御系统;
\item 若处于 \texttt{auto} 模式,系统先检测图像是否疑似对抗样本;
\item 若检测结果超过阈值,则调用 TTC 对图像进行防御修正;
\item 防御后的图像再进入 CLIP 视觉编码器,随后执行动态提示分类。
\end{enumerate}

当前实现提供三种模式。

\begin{enumerate}
\item \texttt{normal}: 不启用防御,按原始动态提示模型推理;
\item \texttt{defend}: 对所有输入统一执行防御;
\item \texttt{auto}: 先检测,仅对疑似对抗样本应用防御。
\end{enumerate}

其中 \texttt{auto} 模式更符合实际使用习惯,因为它试图在鲁棒性和推理开销之间取得平衡。项目文档也明确指出,防御会带来额外计算代价,并且可能轻微影响干净样本精度,因此是否始终启用防御需要结合具体部署需求决定。

\section{安全评估方案}

项目已经提供了独立评估脚本 \texttt{evaluate\_security.py},用于衡量防御机制在干净样本和对抗样本上的影响。该脚本的评估逻辑分为四个阶段:

\begin{enumerate}
\item 关闭防御,评估干净样本准确率;
\item 关闭防御,在 PGD 或 FGSM 攻击下评估对抗样本准确率;
\item 打开防御,评估干净样本准确率;
\item 打开防御,在相同攻击下评估对抗样本准确率。
\end{enumerate}

脚本最终输出四类指标: 无防御干净准确率、有防御干净准确率、无防御对抗准确率、有防御对抗准确率,以及由此计算得到的对抗精度提升和干净样本精度变化。就论文写作而言,这套评估流程已经能够覆盖安全功能是否有效的核心问题。

不过,在当前本地工作区中,安全评估尚未形成可直接写入论文的最终实验数值。原因有两点。第一,当前工作区缺少完整的测试数据目录,无法在本地直接跑完安全评估。第二,现有 \texttt{evaluate\_security.py} 在本地环境中还暴露出数据集类初始化问题,需要继续修正后才能稳定复现。因此,本文在这一章只对安全模块的设计与评估方案进行说明,而不直接给出未经复核的防御实验表格。

\section{本章小结}

本章介绍了项目中新增的系统安全设计。该扩展以 TTC 思路为基础,通过 \texttt{TestTimeCounterattack}、\texttt{AdversarialDetector} 和 \texttt{RobustCustomCLIP} 三部分实现了推理时检测与防御功能。与原有动态提示分类流程相比,安全模块不改变主训练框架,而是作为推理阶段的附加保护机制存在。虽然当前本地环境尚未完成最终安全实验的量化复现,但从代码结构和评估脚本来看,项目已经具备独立开展安全测试的基础。
